


\subsection{Changes in Teachers' Behaviors and School Practices}\label{sec:grading}

\paragraph{Teacher Grading.} Teachers can decide who obtains a preferential seat through their grading. If in response to the percent plan policy they manipulate their grading in a way that weakens the link between achievement and GPA, students in treated schools would have a lower incentive to study to improve their grades. This could explain the negative impacts on effort.\par 
The evidence does not support this mechanism. As shown, pre-college effort reductions resulted in grade reductions (Table \ref{tabreducedform}). Accordingly, the mapping between standardized achievement and grades does not differ between treated and control schools (Table \ref{tabgrading}), suggesting that grading did not respond to the treatment. Consistent with this result, school principals report similar grading practices across treatment groups (Table \ref{tabprincipal}).\par 

\input{Revision/Tables R/tab_grading}

\input{Revision/Tables R/tab_principal}


\paragraph{Teacher Effort and Focus of Instruction.} Teachers could change their focus of instruction (i.e., what portion of the ability distribution they target with their teaching), or they could change effort (class preparation hours and absence days) as an effect of percent plans like PACE. Section \ref{sec:teachervar} describes how we measured these teacher behaviors, and Table \ref{tab:teachereff} shows that there is no evidence that such behaviors responded to the policy.

\input{Revision/Tables R/teachers_effort.tex}




 \paragraph{Schools.} The curriculum is not a possible margin of policy response because the Ministry of Education mandates it. But school principals in treated schools may decide to offer fewer support classes, especially in regards to entrance exam preparation, as performing well on the exam is less critical for an admission. This, in turn, could directly lower students' pre-college achievement, especially in the subjects tested on the exam.\par 
 Using our survey of school principals, we find that treated schools do not differ from control schools regarding the support offered to students (PSU entrance exam preparation support or remedial classes), as shown in Table \ref{tabprincipal}. \par
 Principals may also choose to change the assignment of students to classrooms. We asked them a set of questions on classroom formation, and found no effects, as shown in Table \ref{tabprincipal2}.\par 

 \input{Revision/Tables R/tab_principal2}

    	\subsubsection{Construction of Teacher Variables}\label{sec:teachervar}
		This Section explains how we constructed the teacher variables that enter Table \ref{tab:teachereff} from the survey data that we collected among the Mathematics and Language teachers of the students in our sample.  \par
  
		
		\paragraph{Teacher effort.} For each teacher we observe the hours the teacher spends to prepare his/her classes, and the number of days the teacher was absent from school.\par
		
		\paragraph{Teacher's focus of instruction.} This variable measures whether the teacher is targeting his/her teaching to a specific part of the student ability distribution.\par
	
		For Mathematics and Language teachers separately we  construct a variable indicating the difficulty level at which the teacher is teaching using survey questions about how much of various components of the curriculum the teacher covered during the term, coupled with the teacher's assessment of the difficulty level of each component. For example, for Mathematics we present the teacher with a list of the 4 subfields taken from the official national curriculum (``Algebra and Functions", ``Geometry", ``Statistics and Probability", ``Trigonometry"), and for each subfield we present the teacher with a list of topics taken from the official national curriculum (for example, for ``Algebra and Functions" two topics are ``logarithmic and exponential function and analysis of their graphs" and ``solution of second degree equations"). In all, we presented Mathematics teachers with 13 topics and Language teachers with 11 topics. For each topic, we first ask the teacher what percentage he/she was able to cover during the first semester (which was over when the data collection started). Second, we ask the teacher to think of the average student in his/her $12^{th}$ grade class, and tell us whether he/she thinks that this student would find the topic easy or difficult to understand. The answers to these questions were collected as 5-point Likert scales. Finally, we multiply the coverage and difficulty within each Mathematics (Language) topic and sum over all topics. 
	

\subsection{Reduction in Perceived Returns to College}\label{sec:perceived_returns}


We elicited beliefs about the monetary returns to a college degree at age 30, and about students' awareness of tuition costs. We find that the policy had no impact on students' beliefs about the monetary returns to college (Section \ref{sec:additional_belief_data}), which are large at $82\%$ of age 30 earnings. Such large perceived returns are similar to those measured in \citealp*{hastings2016informed} among Chilean students. The policy had no effect on students' awareness of financial aid ($83.3\%$ of surveyed students are aware they are eligible for a tuition fee waiver, and there is no statistically significant difference between the treatment and control groups (p=0.351)). Therefore, the treatment did not affect students' perceived net returns to college.

    	\subsubsection{Beliefs over Returns to College Degree}\label{sec:additional_belief_data}
		Our survey included the survey instruments developed in \citealp{attanasio2014education} to elicit students beliefs about returns to a college degree. We elicited beliefs about the distribution of wages at age 30 with and without a college degree. We find that students think that, on average, the return to a college degree is 82 percent. This is in line with observed differences in wages between Chileans with and without a college degrees, and in line with results from other surveys on different samples of Chilean high-school students (\citealp*{hastings2016informed}).\par 
		
		We found that the treatment did not have any impact on student beliefs about returns to education (no impacts on the mean nor on the variance of the returns), as reported in Table \ref{expearnings}.
		
		\input{output/tables/exp_earnings}

		Expected (mean) earnings were directly elicited, and we also estimated them, together with the variance of earnings, from elicited c.d.f. values \citep{miglino_tincani_subjective_earnings_2026}. We report results on both measures of expected earnings, for comparison.\par
		The survey questions asked “How much do you expect to earn per month with (without) a college degree on average?” and “How likely are you to earn at least X pesos per month with (without) a college degree?” where X={$200,000$, $800,000$} without a degree and  X={$300,000$, $1,000,000$} with a degree. To calculate the mean and variance of expected earnings using the answers to these questions, we fit the reported c.d.f. values using log-normal distributions for each respondent in the sample. In the estimation sample we kept only the students that answered at least two questions for each scenario (with and without a degree), because we needed at least two c.d.f. values to estimate the mean and variance of the Log-normal distribution. Finally, we used the Generalized Method of Moments to find the mean and variance of the log-normal distribution that minimize the distance of the simulated mean and simulated c.d.f. values from their data analogues.\par
		For variance regressions we use median regressions because the variance is very vulnerable to outlier survey responses in which a student gives the same probability to the likelihood that his/her earnings at age $30$ will be above two different values. 
	

